\chapter{Atrial Flutter}
\index{Atrial Flutter}
\label{sec:Atrial Flutter}

\section{Overview}
Atrial flutter is a supraventricular arrhythmia characterized by a rapid, regular atrial rate of 250--350 beats per minute with a characteristic sawtooth pattern on electrocardiogram. It arises from a macro-reentrant circuit, most commonly in the right atrium around the tricuspid annulus (cavotricuspid isthmus-dependent flutter). This condition shares risk factors with atrial fibrillation and often coexists with or alternates between the two arrhythmias.
\section{Classification}

\begin{description}[font=\normalfont\bfseries,leftmargin=3em,labelwidth=2.5em]
  \item[Cause] Reentrant electrical circuit in the atria
  \item[Organ System] Heart
  \item[Pathophysiology] Macro-reentrant circuit, typically cavotricuspid isthmus-dependent, producing continuous atrial depolarization
  \item[Duration] Acute or Chronic (paroxysmal, persistent, or permanent)
  \item[Onset] Sudden
  \item[Mode of Transmission] Non-communicable
  \item[Clinical Course] Can be paroxysmal, persistent, or permanent; often progresses to atrial fibrillation
  \item[Severity] Mild to Moderate
  \item[Extent of Disease] Localized (right atrial conduction system)
  \item[Age Group] Older adults (peak prevalence after age 60)
  \item[Epidemiology] Prevalence approximately 0.1\% in the general population; increases sharply with age
  \item[ICD-11 Code] BC93.0
\end{description}

\section{Causes}
Atrial flutter is caused by a macro-reentrant circuit within the atrial myocardium, most commonly involving the cavotricuspid isthmus in the right atrium. Conditions that increase atrial pressure or cause atrial enlargement, such as hypertension, heart failure, valvular heart disease, and pulmonary disease, predispose to the development of flutter. Post-cardiac surgery and hyperthyroidism are also well-established triggers.

\section{Risk Factors}

\begin{itemize}[noitemsep]
  \item Advanced age
  \item Hypertension with left atrial enlargement
  \item Heart failure with reduced or preserved ejection fraction
  \item Chronic obstructive pulmonary disease (COPD)
  \item Hyperthyroidism
  \item Obesity and obstructive sleep apnea
  \item Prior cardiac surgery (especially valve or coronary bypass surgery)
\end{itemize}

\section{Signs and Symptoms}

\subsection{Early symptoms}
\begin{itemize}[noitemsep]
  \item Early manifestations of atrial flutter may be subtle and nonspecific
\end{itemize}

\subsection{Common symptoms}
\begin{itemize}[noitemsep]
  \item \subsection{Early symptoms}
  \item \subsection{Common symptoms}
  \item \textbf{Common symptoms:}
  \item Palpitations described as a rapid, regular pounding in the chest
  \item Shortness of breath, especially with exertion
  \item Fatigue and reduced exercise tolerance
  \item Dizziness or lightheadedness
  \item Syncope or near-syncope with rapid ventricular rates
  \item Chest discomfort or pressure
  \item Anxiety or sensation of unease
  \item Pain or discomfort in the affected area
  \item Difficulty performing daily activities
\end{itemize}

\subsection{Less common symptoms}
\begin{itemize}[noitemsep]
  \item Atypical or less common presentations of atrial flutter may occur
\end{itemize}

\subsection{Emergency warning signs}
\begin{itemize}[noitemsep]
  \item Severe or worsening symptoms of atrial flutter require immediate medical evaluation
\end{itemize}
\section{Progression and Stages}
Atrial flutter often presents as paroxysmal episodes that may spontaneously convert to sinus rhythm or progress to persistent flutter requiring cardioversion. Over time, many patients develop concurrent atrial fibrillation (flutter-fib syndrome) as atrial remodeling progresses. Chronic atrial flutter with poorly controlled ventricular rates can lead to tachycardia-induced cardiomyopathy and worsening heart failure.

\section{Complications}

\begin{itemize}[noitemsep]
  \item Thromboembolic events including ischemic stroke (similar risk profile to atrial fibrillation)
  \item Tachycardia-induced cardiomyopathy from persistently rapid ventricular rates
  \item Progression to atrial fibrillation (flutter-fib)
  \item Reduced quality of life
  \item Emotional and psychological effects
\end{itemize}

\section{Diagnosis}

\begin{itemize}[noitemsep]
  \item 12-lead ECG showing characteristic sawtooth flutter waves (typically negative in II, III, aVF)
  \item Holter or event monitoring for paroxysmal episodes
  \item Echocardiogram to assess atrial size and structural heart disease
  \item Laboratory tests including thyroid function and electrolytes
  \item Transesophageal echocardiogram (TEE) to exclude left atrial thrombus before cardioversion
  \item Imaging tests as needed
\end{itemize}

\section{Prevention}

\begin{itemize}[noitemsep]
  \item Treat underlying conditions such as hypertension, heart failure, and hyperthyroidism
  \item Maintain a healthy weight and manage sleep apnea
  \item Avoid excessive alcohol consumption
  \item Anticoagulation therapy based on CHA2DS2-VASc score to prevent stroke
  \item Catheter ablation as preventive curative treatment
  \item Go for regular check-ups
  \item Follow your doctor's screening recommendations
\end{itemize}

\section{Treatment Options}

\begin{itemize}[noitemsep]
  \item Rate control with beta-blockers or calcium channel blockers
  \item Electrical or pharmacologic cardioversion to restore sinus rhythm
  \item Catheter ablation (cavotricuspid isthmus ablation) with high success rate
  \item Anticoagulation (warfarin or direct oral anticoagulants) for stroke prevention
  \item Antiarrhythmic drugs (flecainide, propafenone, amiodarone) for rhythm control
  \item Lifestyle changes and self-care
  \item Surgery when needed (maze procedure in selected cases)
\end{itemize}

\section{Living with It}

\begin{itemize}[noitemsep]
  \item Follow your treatment plan as prescribed
  \item Keep up with regular appointments
  \item Adjust your daily habits as needed
  \item Reach out to family, friends, or support groups
  \item Keep learning about your condition
\end{itemize}

\section{Outlook}
The prognosis for patients with atrial flutter is generally favorable with appropriate management. Catheter ablation of typical atrial flutter has a success rate of 90--95\% and is considered first-line definitive therapy in suitable candidates.
